\hojaejercicio{Aplicaciones Teorema de Cauchy II }
\ejercicio{Demuestra la fórmula integral de Poisson: Si $f$ es holomorfa en un 
dominio $\Omega \subset \mathbb{C}$ tal que $\Omega \supset \overline{D}(0;1)$, 
entonces para cada $z=re^{i\theta} \in D(0;1)$ se tiene:

$$ f(re^{i\theta}) = \frac{1}{2\pi} \int_{0}^{2\pi} \frac{1-r^2}{1-2r\cos(\theta-t)+r^2} \cdot f(e^{it}) \, dt $$

(ind.: aplica la fórmula integral de Cauchy a $f(z)g_z(z)$, siendo $g_z(w)=(1-r^2)/(1-w\bar{z}))$.}{
    Sea $f \in H(\Omega)$ y $\overline{D}(0,1) \subset \Omega$.

Tenemos $h_z(w) = f(w) \cdot g_z(w) = f(w) \cdot \frac{1-r^2}{1-w\bar{z}}$.

$h(z)$ es holomorfa en $\Omega \setminus \{1/\bar{z}\}$ puesto que $f \in H(\Omega)$ y $g \in H(\mathbb{C} \setminus \{1/\bar{z}\})$ al ser cocientes de funciones holomorfas cuyo denominador se anula en $w = \frac{1}{\bar{z}}$ y producto de holomorfas en $\mathbb{C} \setminus \{1/\bar{z}\}$
\[ g_z(w) \cdot f(w) \in H(\Omega \cap (\mathbb{C} \setminus \{1/\bar{z}\}))=H(\Omega \setminus \{1/\bar{z}\})\]

Cogiendo entonces $z = re^{i\theta} \in D(0,1)$, entonces $\bar{z} \in D(0,1)$ (porque $|z|=|\bar{z}|$).
$$ \left| \frac{1}{\bar{z}} \right| = \frac{1}{|\bar{z}|} > 1 \implies \frac{1}{\bar{z}} \notin \overline{D}(0,1) \quad \text{por lo tanto, } \overline{D}(0,1) \subset \mathbb{C} \setminus \{1/\bar{z}\}. $$

Al tener $\overline{D}(0,1) \subset \Omega$ y $\overline{D}(0,1) \subset \mathbb{C} \setminus \{1/\bar{z}\}$, entonces:
$$ \overline{D}(0,1) \subset \Omega \cap \mathbb{C} \setminus \{1/\bar{z}\} = \Omega \setminus \{1/\bar{z}\} $$
que es abierto conexo en $\mathbb{C}$.

$h \in H(\Omega \setminus \{1/\bar{z}\})$, así pues, podemos aplicar la fórmula integral de Cauchy.

\begin{align*}
h(z) &= \frac{1}{2\pi i} \int_{\partial D(0,1)} \frac{h(w)}{w-z} \, dw \\
&= \frac{1}{2\pi i} \int_{C(0,1)} \frac{f(w) \cdot g_z(w)}{w-z} \, dw \\
&= \frac{1}{2\pi i} \int_{C(0,1)} \frac{f(w) \frac{1-r^2}{1-w\bar{z}}}{w-z} \\
&= \frac{1}{2\pi i} \int_{C(0,1)} f(w) \frac{1-r^2}{(w-z)(1-w\bar{z})} \, dw
\end{align*}

Como $w \in C(0,1) \iff |w|=1 \iff |w|^2=1 \iff w\bar{w}=1$:

\begin{align*}
\frac{1}{2\pi i} \int_{C(0,1)} f(w) \frac{1-r^2}{(w-z)(\bar{w}-\bar{z})w} \, dw 
&= \frac{1}{2\pi i} \int_{C(0,1)} f(w) \frac{1-r^2}{|w-z|^2 w} \\
&= \frac{1}{2\pi} \int_{0}^{2\pi} f(e^{it}) \frac{(1-r^2)  e^{it}}{|e^{it}-re^{i\theta}|^2  e^{it}} \, dt \\
&= \frac{1}{2\pi} \int_{0}^{2\pi} f(e^{it}) \frac{(1-r^2)}{|e^{it}(1-re^{i\theta})|^2} \, dt \\
&= \frac{1}{2\pi} \int_{0}^{2\pi} f(e^{it}) \frac{1-r^2}{|1-re^{i(\theta-t)}|^2} \, dt \quad (*)
\end{align*}

Desarrollando concluimos:
\begin{align*}
|1 - re^{i(\theta-t)}|^2 &= (1 - re^{i(\theta-t)})(1 - re^{-i(\theta-t)}) \\
&= 1 - re^{-i(\theta-t)} - re^{i(\theta-t)} + r^2 \\
&= 1 + r^2 - 2r\cos(\theta-t)
\end{align*}

Sustituyendo en $(*)$:
$$ (*) = \frac{1}{2\pi} \int_{0}^{2\pi} f(e^{it}) \frac{1-r^2}{1+r^2-2r\cos(\theta-t)} \, dt $$
}

\ejercicio{

Calcula $\displaystyle \int_{C(0;1)} \frac{\text{Re } z}{z - \frac{1}{2}} , dz$ \quad $(\text{recuerda que } \bar{z} = \frac{1}{z} \text{ si } |z|=1)$.
}{

$$ \int_{C(0,1)} \frac{\text{Re } z}{z - \frac{1}{2}} , dz. \quad \text{Sabemos que } 2\text{Re } = z + \overline{z} \Rightarrow \text{Re } = \frac{z+\overline{z}}{2} \Rightarrow $$

\begin{align*}
\int_{C(0,1)} f(z) , dz &= \int_{C(0,1)} \frac{\frac{z+\overline{z}}{2}}{z-\frac{1}{2}} , dz = \int_{C(0,1)} \frac{z+\frac{1}{z}}{2z-1} , dz = \int_{C(0,1)} \frac{z}{2z-1} , dz + \int_{C(0,1)} \frac{1/z}{2z-1} \, dz \\
&= \int_{C(0,1)} \frac{z}{2} \cdot \frac{1}{z-\frac{1}{2}} , dz + \int_{C(0,1)} \frac{1}{2z} \cdot \frac{1}{z-\frac{1}{2}} , dz.
\end{align*}

entonces: $\Rightarrow \int_{C(0,1)} \frac{z}{2} \cdot \frac{1}{z-1/2} \, dz = \frac{2\pi i \cdot (1/2)}{2} \quad$ por la fórmula integral de Cauchy.

Si cogemos dos discos de centro $0$ y $1/2$ con radio $\delta$ tq $C(0,\delta) \cap C(1/2, \delta) = \emptyset$.
Entonces, cogiendo la parametrización adecuada de la curva que une las tres circunferencias (y viendo que es homótopa a 0). Nos resulta en:

\begin{align*}
\int_{C(0,1)} \frac{1}{2z} \cdot \frac{1}{z-1/2} , dz &= -\int_{C(0,\delta)} \frac{1}{z} \cdot \frac{1}{2z-1} , dz - \int_{C(1/2,\delta)} \frac{1}{2z} \cdot \frac{1}{z-1/2} , dz \\
&= -2\pi i \left(\frac{1}{2\cdot 0 - 1}\right) - 2\pi i \left(\frac{1}{2 \cdot \frac{1}{2}}\right) \\
&= -2\pi i + 2\pi i = \boxed{0}
\end{align*}

\[\int_{C(0,1)} \frac{\text{Re } z}{z - 1/2} , dz = \int_{C(0,1)} \frac{z}{2} \frac{1}{z-1/2} , dz = \frac{2\pi i \cdot 1/2}{2} = \boxed{\frac{\pi i}{2}} \]
}

%ejercicio 3
\ejercicio{Sea $\log z$ el logaritmo principal de $z$ en $\Omega = \mathbb{C} \setminus (-\infty, 0]$ y sea $a = \left( \frac{e}{\sqrt{2}} \right)(1+i)$.
Calcula $\displaystyle \int_{C(a;1)} \frac{\log z}{z-a} \, dz$.
}{
    Tenemos que $\log z$ es el logaritmo principal de $z$ en $\Omega = \mathbb{C} \setminus (-\infty, 0]$ por lo que:
$$ \text{Log } z = \log|z| + i\text{Arg}(z) $$
siendo $\text{Arg}(z)$ el argumento principal de $z$ en $\Omega$, i.e., $\text{Arg}(z) \in (-\pi, \pi)$.

Tenemos que es continua en $\Omega$ y además, es holomorfa porque es una determinacion del logaritmo continua en $\Omega$ que es un abierto conexo (prop 1.3.1) (es conexo por caminos).
Por lo tanto $\text{Log}(z) \in H(\Omega)$ con $(\log z)' = \frac{1}{z} \quad \forall z \in \Omega$.

Ahora, teniendo en cuenta que:
\begin{align*}
\text{Re } z \ge -1 \quad \forall z \in \overline{D}(0,1) &\implies \text{Re } z \ge \text{Re } a - 1 \quad \forall z \in \overline{D}(a,1) \\
&\implies \text{Re } z \ge \frac{e}{\sqrt{2}} - 1 > 0 \\
\text{Im } z \ge -1 \quad \forall z \in \overline{D}(0,1) &\implies \text{Im } z \ge \text{Im } a - 1 \quad \forall z \in \overline{D}(a,1) \\
&\implies \text{Im } z \ge \frac{e}{\sqrt{2}} - 1 > 0
\end{align*}
$\implies \overline{D}(a,1) \subset \Omega$.

Entonces utilizando la fórmula integral de Cauchy:

\begin{align*}
\int_{C(a,1)} \frac{\log z}{z-a} \, dz &= 2\pi i \log(a) \\
&= 2\pi i \left( \log e + \frac{\pi}{4}i \right) \\
&= \boxed{2\pi i  - \frac{\pi^2}{2}}
\end{align*}
}

%ejercicio 4
\ejercicio{
    Utiliza el principio de identidad para probar que
\[sen(z+w) = sen z cos w + cos z sen w \quad \forall z,w \in \mathbb{C} \]
}
{
    El principio de identidad nos dice que si $\Omega$ es un dominio y $g, f \in H(\Omega)$ entonces si el conjunto de puntos donde son iguales tiene algún punto de acumulación entonces $f \equiv g$ en $\Omega$.

Por lo que si consideramos $\Omega = \mathbb{C}$ y fijamos $z \in \mathbb{R}$.

Entonces:
\[sen(z+w) = sen z cos w + cos z sen w \quad \forall w \in \mathbb{R} \]

Como $\mathbb{R}$ tiene puntos de acumulación en $\Omega = \mathbb{C}$. Entonces se da que para dicho $z \in \mathbb{R}$:
\[ sen(z+w) = sen z cos w + cos z \cdot sen w \quad \forall w \in \mathbb{C}. \]

Entonces fijamos ahora $w \in \mathbb{C}$. Se da que:
\[ sen(z+w) = sen z cos w + cos z \cdot sen w \quad \forall z \in \mathbb{R} \]
por lo anterior, entonces, por el principio de identidad:
\[ sen(z+w) = sen z cos w + cos z sen w \quad \forall z \in \mathbb{C}. \]

Es decir:
\[sen(z+w) = sen z cos w + cos z sen w \quad \forall z, w \in \mathbb{C}. \]
}
%ejercicio 5
\ejercicio{
    Sea $D = D(0;1)$. Para cada $a \in D$ consideramos la transformación de Möbius:
    \[
    \varphi_a(z) = \frac{z - a}{1 - \bar{a}z}, \quad \forall z \in \bar{D}
    \]
    Demuestra que $\varphi_a(D) = D$, $\varphi_a(C(0,1)) = C(0,1)$, $\varphi_a: D \to D$ es biyectiva, $\varphi_a(0) = -a$, $\varphi_a(a) = 0$, $(\varphi_a)^{-1} = \varphi_{-a}$, $\varphi_a'(0) = 1 - |a|^2$ y $\varphi_a'(a) = \frac{1}{1 - |a|^2}$.
}{
    Necesitamos demostrar que $\varphi_{a}(D)=D$. O lo que es lo mismo:
\[
\text{Si } |z|<1 \implies |\varphi_{a}(z)|<1.
\]
Recordamos la definición: $\varphi_{a}(z) = \frac{z-a}{1-\overline{a}z}$.

\subsection*{1. Demostración de la imagen del disco}
Entonces, calculamos el módulo al cuadrado:
\[
|\varphi_{a}(z)|^2 = \left| \frac{z-a}{1-\overline{a}z} \right|^2 = \frac{|z-a|^2}{|1-\overline{a}z|^2} = \frac{(z-a)(\overline{z}-\overline{a})}{(1-\overline{a}z)(1-a\overline{z})}
\]
Queremos ver si esto es $< 1$. Desarrollamos la desigualdad:
\[
\frac{|z|^2 + |a|^2 - \overline{a}z - a\overline{z}}{1 + |a|^2|z|^2 - \overline{a}z - a\overline{z}} < 1
\]
Multiplicando por el denominador (que es positivo):
\[
|z|^2 + |a|^2 - \overline{a}z - a\overline{z} < 1 + |a|^2|z|^2 - \overline{a}z - a\overline{z}
\]
Simplificamos cancelando los términos comunes ($-\overline{a}z - a\overline{z}$):
\[
|z|^2 + |a|^2 < 1 + |a|^2|z|^2
\]
Reordenamos los términos:
\[
|z|^2 + |a|^2 - |a|^2|z|^2 - 1 < 0
\]
Factorizamos:
\[
|z|^2(1-|a|^2) - (1-|a|^2) < 0 \iff (1-|a|^2)(|z|^2 - 1) < 0
\]
O equivalentemente:
\[
|z|^2 + |a|^2(1-|z|^2) < 1 \iff \frac{1-|a|^2}{1-|a|^2|z|^2} \dots \quad (\text{Forma simplificada: } (1-|a|^2)(1-|z|^2) > 0)
\]
Como $|a|<1$ y $|z|<1$, la desigualdad se cumple. Por lo tanto, $|\varphi_{a}(z)| < 1$.


Sea $z_0 \in D$, buscamos $z \in D$ tal que $\varphi_{a}(z) = z_0$.
\[
\varphi_{a}(z) = z_0 \iff \frac{z-a}{1-\overline{a}z} = z_0 \iff z-a = z_0(1-\overline{a}z)
\]
\[
z - a = z_0 - z_0\overline{a}z \iff z + z z_0\overline{a} = z_0 + a \iff z(1+z_0\overline{a}) = z_0 + a
\]
\[
z = \frac{z_0 + a}{1 + z_0\overline{a}}
\]
¿Está este $z$ en $D(0,1)$?
\[
\left| \frac{z_0+a}{1+\overline{a}z_0} \right| < 1
\]
(El desarrollo es análogo al apartado anterior intercambiando los roles de $z$ y $a$).
Por lo tanto, $\varphi_a(D) = D$.



\[
\varphi_{a}(C(0,1)) = C(0,1)
\]
Queremos ver que: Si $|z|=1 \implies |\varphi_{a}(z)|=1$ y además que todo punto de $C(0,1)$ es imagen de otro punto de $C(0,1)$.
Demostración:
\[
\text{Si } |z|=1 \implies |\varphi_{a}(z)| = \left| \frac{z-a}{1-\overline{a}z} \right| = 1 \iff |z-a| = |1-\overline{a}z|
\]
Elevando al cuadrado:
\[
|z|^2 + |a|^2 - a\overline{z} - \overline{a}z = 1 + |a|^2|z|^2 - a\overline{z} - \overline{a}z
\]
Como $|z|=1$, entonces $|z|^2=1$:
\[
1 + |a|^2 = 1 + |a|^2(1)
\]
La igualdad se cumple.
La sobreyectividad se demuestra de forma análoga al apartado anterior.

$\varphi_{a}: D \to D$ es biyectiva. La sobreyectividad es la hemos demostrado anterioremente, falta comprobar inyectividad.
Si $\varphi_{a}(w) = \varphi_{a}(z)$:
\[
\frac{w-a}{1-w\overline{a}} = \frac{z-a}{1-z\overline{a}} \iff (w-a)(1-z\overline{a}) = (z-a)(1-w\overline{a})
\]
\[
w - wz\overline{a} - a + az\overline{a} = z - zw\overline{a} - a + aw\overline{a}
\]
Simplificamos términos iguales:
\[
w + z|a|^2 = z + w|a|^2 \iff w - z = w|a|^2 - z|a|^2
\]
\[
(w-z) = |a|^2(w-z) \iff (1-|a|^2)(w-z) = 0
\]
Como $|a|<1$, entonces $1-|a|^2 \neq 0$, por lo que $w-z=0 \implies w=z$.

Para el resto, son simples comprobaciones:
\[
\varphi_{a}(a) = \frac{a-a}{1-|a|^2} = 0
\]
\[
\varphi_{a}(0) = \frac{0-a}{1-0} = -a
\]
El resto se demuestran de forma trivial.
}

% Ejercicio 6
\ejercicio{
    Utiliza el lema de Schwarz y el ejercicio anterior para probar que si $f: D \to D$ es biholomorfa (= biyectiva, holomorfa con inversa holomorfa), entonces $f$ es una transformación de Möbius. Deduce que si $D_1$ y $D_2$ son discos y $g: D_1 \to D_2$ es biholomorfa, entonces $g$ es una transformación de Möbius.
}{
  $f: D \to D$ biholomorfa.

\vspace{0.5cm}

Debemos usar el Lema de Schwarz: Sea $D=D(0,1)$ y $f \in H(D)$ tal que $|f(z)| \le 1 \quad \forall z \in D$ y $f(0)=0$. Entonces:
\[
|f(z)| \le |z| \quad \forall z \in D \quad \text{y} \quad |f'(0)| \le 1
\]
Además, si para algún $z_0 \in D \setminus \{0\}$ se cumple que $|f(z_0)| = |z_0|$ o si $|f'(0)|=1$, entonces $f$ es una rotación.
Sabemos que $f^{-1}: D \to D$ es holomorfa.

\vspace{0.3cm}

Sin embargo, nuestra función $f$ no tiene por qué cumplir que $f(0)=0$. Por lo tanto:

Construyamos una función que reúna las condiciones. Nuestra $f$ es biyectiva, así que $\exists w \in D$ tal que $f(w)=0$. Por lo tanto, usando $\varphi_{-w}(z) = \frac{z+w}{1+\bar{w}z}$, tenemos entonces que:
\[
0 \xrightarrow{\varphi_w} w \xrightarrow{f} 0
\]
\[
\varphi_{-w}(f(0))=0
\]

Sea $h(z) = \varphi_{-w} \circ f(z)$.

\vspace{0.3cm}

Al ser $f$ y $\varphi_{-w}$ biyectivas en $D$, $h$ es biyectiva.

$f \in H(D)$ y $\varphi_{-w} \in H(\mathbb{C} \setminus \{ 1/\bar{w} \})$. Pero como $w \in D(0,1)$, entonces $|1/\bar{w}| = 1/|w| > 1$. Por lo tanto $D \subset \mathbb{C} \setminus \{ 1/\bar{w} \}$, lo que implica que $\varphi_{-w} \in H(D)$.

Así, $h = \varphi_{-w} \circ f \in H(D)$ al ser composición de holomorfas. Además $h(D) \subseteq D$.
\[
\implies h(0)=0 \quad \text{y} \quad |h(z)| \le 1 \quad \forall z \in D
\]
Por lo que podemos aplicar el lema. Concluimos que:
\[
\begin{cases} 
|h(z)| \le |z| \quad \forall z \in D \\ 
|h'(0)| \le 1 
\end{cases}
\]

\vspace{0.3cm}

Como $\varphi_{-w}$ y $f$ son biholomorfas en $D \implies h = \varphi_{-w} \circ f$ es biholomorfa en $D$.
$\implies h^{-1}$ es holomorfa en $D$ y cumple las condiciones del Lema de Schwarz.

Trivialmente: $h^{-1}(h(z)) = z$. Derivando por la regla de la cadena:
\[
(h^{-1})'(h(z)) \cdot h'(z) = 1 \implies (h^{-1})'(h(0)) = \frac{1}{h'(0)}
\]
Como $h(0)=0$:
\[
h^{-1}(0) = 0 \implies (h^{-1})'(0) = \frac{1}{h'(0)}
\]

Al cumplir $h^{-1}$ el Lema de Schwarz, sabemos que $|(h^{-1})'(0)| \le 1$. Sustituyendo:
\[
\left| \frac{1}{h'(0)} \right| \le 1 \implies |h'(0)| \ge 1
\]

Llegamos a que $|h'(0)| \ge 1$ y anteriormente teníamos $|h'(0)| \le 1$.
\[
\implies |h'(0)| = 1
\]

Entonces, por la parte de la igualdad del propio lema, $h: D \to D$ es una \textbf{rotación} tal que:
\[
h(z) = e^{i\theta}z \quad \text{con } \theta \in [0, 2\pi)
\]

Como $h = \varphi_{-w} \circ f$, deshacemos el cambio:
\[
e^{i\theta}z = \varphi_{-w}(f(z)) \implies (\varphi_{-w})^{-1}(e^{i\theta}z) = f(z)
\]
\[
\implies f(z) = \varphi_{w}(e^{i\theta}z)
\]

$f$ es composición de Transformaciones de Möbius (rotación y automorfismo), por lo tanto \textbf{$f$ es una Transformación de Möbius}.
}
\ejercicio{Sea $f(z) = z + \frac{1}{z}$. Prueba que $f$ es inyectiva en $\mathbb{C} \setminus \bar{D}(0;1)$. Determina $f(C(0,r))$ para $r > 1$ y para $r = 1$.}{}
\ejercicio{Sea $f: D(0;1) \to \mathbb{C}$ holomorfa y acotada con $f(0) = 0$. Usa el lema de Schwarz para probar que $\sum_{n=1}^{\infty} f(z^n)$ converge en $D = D(0;1)$ y su función suma es holomorfa en $D$.}{
    Tenemos que $f \in H(D)$ y $|f(z)| \le M$ para todo $z \in D(0,1)$.
Consideramos la función:
\[
g(z) = \frac{f(z)}{M}
\]
Esta función cumple las premisas del Lema de Schwarz ($g(0)=0$ y $|g(z)| \le 1$).
\[
\implies \left| \frac{f(z)}{M} \right| \le |z| \quad \text{y} \quad \left| \frac{f'(0)}{M} \right| \le 1
\]
\[
\implies |f(z)| \le M|z|
\]
Aplicando esto a los términos de la serie:
\[
|f(z^n)| \le M|z^n| = M|z|^n
\]
Por el criterio de comparación de series, estudiamos:
\[
\sum_{n=1}^{\infty} |f(z^n)| \le \sum_{n=1}^{\infty} M|z|^n
\]
Como $|z|<1$, la serie geométrica $\sum |z|^n$ converge. Por lo tanto:
\[
\sum_{n=1}^{\infty} f(z^n) \quad \text{converge absolutamente para } z \in D(0,1).
\]

\subsubsection*{Convergencia uniforme sobre compactos}
Sea $f_n(z) = f(z^n)$. Veamos que la serie converge uniformemente sobre compactos de $D(0,1)$.
Sea $K \subset D(0,1)$ un compacto.
Como $K$ es cerrado y acotado y $K \subset D(0,1)$, existe una distancia al borde $>0$. La función $|z|$ alcanza un máximo en $K$.
Sea $r = \max \{|z| : z \in K \} < 1$.
\[
|f_n(z)| \le M|z|^n \le M r^n \quad \forall z \in K
\]
Como $\sum M r^n$ converge (serie geométrica con $r<1$), por el \textbf{Criterio M de Weierstrass}, la serie $\sum f_n(z)$ converge uniformemente sobre compactos de $D$.

Sea $g_m(z) = \sum_{n=1}^{m} f(z^n)$.
Como $g_m \in H(D)$ y la sucesión converge uniformemente sobre compactos, por el Teorema de Weierstrass, el límite $g(z)$ es una función holomorfa.
\[
g(z) = \sum_{n=1}^{\infty} f(z^n) \in H(D)
\]
Además, la derivada se puede calcular término a término:
\[
g'(z) = \sum_{n=1}^{\infty} (f(z^n))' = \sum_{n=1}^{\infty} f'(z^n) \cdot n z^{n-1}
\]}